\subsection{Spectral gap used in the CI proof}
\label{sec:fixed-girth-ci}\label{app:fixed-girth-ci}

This subsection proves \eqref{eq:girth5-poincare} on every residual instance
satisfying \eqref{eq:girth5-degree-budget} and having girth at least five.
We use $\Deg_0(\delta)$ and $\gamma_\delta$ from
\eqref{eq:girth5-gap-constants}, under the same hypothesis
$q\ge(1+\delta)\Deg$.  Throughout, $P_v$
is conditional expectation given all spins except $\sigma_v$, and
$Q_v:=I-P_v$.  Heat-bath updates occur at rate one per free vertex.
All operator inequalities are on the supported $L^2$ space of the current
Gibbs law: $A\succeq B$ means that $A-B$ is positive semidefinite.
We write $\|f\|_\mu^2:=\mathbb E_\mu|f|^2$ and similarly for $\rho$.
Unsubscripted function norms are in $L^2(\mu)$, and unsubscripted operator
norms are the induced Hilbert-space norms.  Finite colour matrices instead
carry the subscript $2\to2$.  The identity $I$ acts on the space in
question, and $\{A,B\}:=AB+BA$ denotes the anticommutator.

\begin{lemma}[Star bounds and singleton complement]
\label{lem:girth5-star-bounds}
Fix $0<\delta\leq1$, $\Deg\geq\Deg_0(\delta)$,
$q\ge(1+\delta)\Deg$, and $x\in[0,1]$, and put
$s=1-x$.  Consider a supported conditional star from a residual instance
satisfying \eqref{eq:girth5-degree-budget} at every free vertex.  Its centre
has $d$ free leaves and $k$ fixed neighbours, with $d+k\leq\Deg$.  Write
\[
 \mu(c,\sigma)=\frac{a_c}{Z}\prod_{i=1}^{d}
 \bigl[p_i(\sigma_i)(1-s\mathbf1_{\{\sigma_i=c\}})\bigr],
 \qquad a_c=x^{b(c)},\quad \sum_cb(c)=k,
\]
where $b(c)$ counts fixed neighbours of colour $c$ and each $p_i$ is
the cavity law obtained after deleting the centre edge.
Here $c\in\colours$ is the centre colour and
$\sigma=(\sigma_1,\ldots,\sigma_d)$ is the leaf configuration.  The
normalizing constant is $Z=\sum_c a_c\prod_i(1-sp_i(c))$.
Let $\nu$ be the centre marginal and $\rho=\bigotimes_i p_i$ the product
cavity law on the leaves; $\mu_{\rm leaves}$ denotes the actual leaf
marginal.  Set
\[
 B_{\rm SG}=\frac{1}{\delta\Deg},\qquad
 \kappa=\frac{B_{\rm SG}^2}{(1-B_{\rm SG})^2},\qquad
 \omega_{\rm SG}=\frac{sB_{\rm SG}}{(1-B_{\rm SG})^2},\qquad c_\delta=\frac{\delta}{1+\delta}.
\]
Then, uniformly over all supported stars (including $x=0$),
\begin{equation}\label{eq:sg-star-density}
 s p_i(c)\le B_{\rm SG},\quad s\nu(c)\le B_{\rm SG},\quad
 \mathcal R:=\frac{d\mu_{\rm leaves}}{d\rho}\ge c_\delta.
\end{equation}
Let $P$ be the centre heat-bath projection and $Q=I-P$, and let
$P_i$ and $Q_i=I-P_i$ denote the heat-bath projection and its complement at
leaf $i$.  Decompose the mean-zero space into the conditional leaf Hoeffding
sectors $K_0\oplus K_1\oplus K_{\ge2}$: $K_0$ consists of mean-zero
functions of the centre, and $K_1$ is the orthogonal sum of the spaces of functions
$g_i(c,\sigma_i)$ with $\mathbb E[g_i\mid c]=0$, and $K_{\ge2}$ is the
sum of higher conditional Hoeffding sectors, involving at least two
distinct conditionally centred leaf factors and allowing coefficients
that depend on $c$.
Write $K_+:=K_1\oplus K_{\ge2}$.  For a closed subspace $\mathcal S$,
$\Pi_{\mathcal S}$ denotes orthogonal projection in $L^2(\mu)$.
Block subscripts indicate the target and source sectors; for example,
$P_{ab}:=\Pi_{K_a}P|_{K_b}$ for $a,b\in\{0,+\}$.  Put
\[
 C=P_{+0}(I-P_{00})^{-1}P_{0+},\qquad H=I-P_{++}-C,
\]
whenever $r:=\|P_{00}\|<1$, and let $H_{11}$ be the compression of $H$
to $K_1$.  If $\mathcal A$ is the subspace generated by
leaf-additive residuals
$h_i(\sigma_i)-\mathbb E[h_i(\sigma_i)\mid c]$, and
$\mathcal B=K_1\ominus\mathcal A$, then
\[
 r\le r_*:=\frac{d\kappa(1+\omega_{\rm SG})^{d-1}}{c_\delta},\qquad
 H_{11}|_{\mathcal B}\succeq(1-\eta)I,
\quad \eta:=r_*+\varepsilon_1,
\]
where
\begin{equation}\label{eq:sg-singleton-complement}
 \varepsilon_1=
 \begin{cases}
 \dfrac{2(d-1)\kappa(1+\omega_{\rm SG})^{d-2}}
       {c_\delta(1-B_{\rm SG})},&d\ge2,\\[3pt]
 0,&d\le1.
 \end{cases}
\end{equation}
Moreover, for $d\ge1$, put $\theta_0=(1+\delta/2)^{-1}$.  Then
\begin{equation}\label{eq:sg-additive-compression}
 \Pi_{\mathcal A}C\Pi_{\mathcal A}\preceq \frac{\theta_0}{d}\Pi_{\mathcal A}.
\end{equation}
\end{lemma}

\begin{proof}
If $d=0$, then $\mathcal A=\{0\}$ and there is no additive sector to
control.  All estimates below containing $1/d$, including
\eqref{eq:sg-additive-compression}, are used only when $d\ge1$.
At a free vertex with $k$ fixed and $d$ free neighbours, let $Z_{\rm loc}$
be the sum of its $q$ conditional colour weights after the other spins are
fixed.  The degree budget gives
$Z_{\rm loc}\ge q-s(k+d)\ge\delta\Deg$.  Hence every one-site
marginal satisfies $sp_i(c),s\nu(c)\le B_{\rm SG}$.  If
$A_0=\sum_ca_c$, then $A_0\ge q-sk\ge s(d+\delta\Deg)$ and
\[
 \mathcal R(\sigma)=\frac{\sum_ca_c\prod_i(1-s\mathbf1_{\{\sigma_i=c\}})}
 {\sum_ca_c\prod_i(1-sp_i(c))}
 \ge \frac{A_0-sd}{A_0}
 \ge \frac{\delta\Deg}{d+\delta\Deg}\ge c_\delta .
\]
At $s=0$ this is interpreted by continuity.  At $x=0$, centre-colour
coordinates and sums weighted by $\nu$ may be restricted to
$\mathcal C_*=\{c:\nu(c)>0\}$; zero-$\nu$ terms vanish.
Each leaf expectation and sum retains the full support of its own
cavity law $p_i$.

For the complement estimate define
\[
 l_i^c=\frac{1-s\mathbf1_{\{\sigma_i=c\}}}{1-sp_i(c)},
 \qquad \xi_i^c=l_i^c-1,
 \qquad \mathcal R=\sum_c\nu(c)\prod_i l_i^c .
\]
The one-leaf Gram matrix is
\[
 \Psi_i(c,c')=\mathbb E_{p_i}[\xi_i^c\xi_i^{c'}]
 =\frac{s^2(\mathbf1_{c=c'}p_i(c)-p_i(c)p_i(c'))}
 {(1-sp_i(c))(1-sp_i(c'))}.
\]
It is positive semidefinite and
\begin{equation}\label{eq:sg-gram-bound}
 \|\operatorname{diag}(\sqrt\nu)\Psi_i\operatorname{diag}(\sqrt\nu)\|_{2\to2}
 \le\kappa,\qquad \max_c\Psi_i(c,c)\le \omega_{\rm SG}.
\end{equation}
The Schur product norm bound for positive semidefinite matrices therefore
implies, for every nonempty $J\subseteq[d]$, the following bound, where
$\circ$ denotes entrywise (Schur) multiplication:
\begin{equation}\label{eq:sg-schur-product}
 \left\|\operatorname{diag}(\sqrt\nu)
 \Big(\mathop{\circ}_{j\in J}\Psi_j\Big)
 \operatorname{diag}(\sqrt\nu)\right\|_{2\to2}
 \le\kappa \omega_{\rm SG}^{|J|-1}.
\end{equation}

If $g=\sum_i g_i\in K_1\ominus\mathcal A$, then orthogonality to every
$h_i(\sigma_i)-\mathbb E[h_i\mid c]$, together with
$\mathbb E[g_i\mid c]=0$, gives
$\mathbb E_\mu[g_i\mid\sigma_i]=0$.  Bayes' formula therefore gives
$\sum_c\nu(c)g_i(c,\sigma_i)l_i^c=0$; this is the vanishing constant
term in the following $L^2(\rho)$-orthogonal expansion (and
$\mathbb E_{p_i}[g_i(c,\cdot)l_i^c]=0$ ensures that the displayed terms
indexed by distinct $I$ are in distinct product Hoeffding sectors):
\begin{equation}\label{eq:sg-hoeffding-expansion}
 (Pg)\mathcal R=\sum_{\substack{I\subseteq[d]\\|I|\ge2}}
 \sum_{i\in I}\sum_c \nu(c)g_i(c,\sigma_i)l_i^c
 \prod_{j\in I\setminus\{i\}}\xi_j^c .
\end{equation}
For fixed $I,i$, \eqref{eq:sg-schur-product}, followed by $(l_i^c)^2\le l_i^c/(1-B_{\rm SG})$, yields
\begin{equation}\label{eq:sg-component-bound}
 \left\|\sum_c\nu(c)g_i(c,\sigma_i)l_i^c
 \prod_{j\in I\setminus\{i\}}\xi_j^c\right\|_\rho^2
 \le \frac{\kappa \omega_{\rm SG}^{|I|-2}}{1-B_{\rm SG}}\|g_i\|_\mu^2.
\end{equation}
Summing the orthogonal Hoeffding components and using
$\|\sum_{i\in I}u_i\|^2\le |I|\sum_i\|u_i\|^2$ gives
\[
 \|(Pg)\mathcal R\|_\rho^2\le \frac{2(d-1)\kappa(1+\omega_{\rm SG})^{d-2}}{1-B_{\rm SG}}\|g\|_\mu^2.
\]
Since $\|Pg\|_\mu^2=\mathbb E_\rho[((Pg)\mathcal R)^2/\mathcal R]$,
\eqref{eq:sg-star-density} gives $\|Pg\|_\mu^2\le\varepsilon_1\|g\|_\mu^2$
with $\varepsilon_1$ as in \eqref{eq:sg-singleton-complement}.
Together with $C\preceq rI$ below, this yields
$H_{11}|_{\mathcal B}\succeq(1-r_*-\varepsilon_1)I$.

For a mean-zero function $h$ of the centre alone, the expansion of
$(Ph)\mathcal R$ starts in degree one;
\eqref{eq:sg-schur-product} and the same calculation give $\|Ph\|^2\le r_*\|h\|^2$,
which is the bound on $r$.  The projection identities
$P_{0+}P_{+0}=P_{00}-P_{00}^2$ and
$P_{++}P_{+0}=P_{+0}(I-P_{00})$ imply $H^2=H$ and
$0\preceq C\preceq rI$.

It remains to prove \eqref{eq:sg-additive-compression}.  For a leaf-additive function
$h=\sum_i h_i$, write $\bar h=\mathbb E[h\mid c]$ and $a=h-\bar h$.
Conditional independence gives
$a\in\mathcal A$ and
$\|a\|^2=\sum_i\mathbb E\operatorname{Var}(h_i\mid c)$.  The exact block
relations for $Ph=h$ imply
\begin{equation}\label{eq:sg-block-identity}
 (I-P_{00})\bar h=P_{0+}a,\qquad
 \langle a,Ca\rangle=\|Q\bar h\|^2.
\end{equation}
After subtracting constants from each $h_i$, normalize
$\mathbb E_{p_i}h_i=0$.  Then
$\bar h_i(c):=\mathbb E[h_i(\sigma_i)\mid c]=-sp_i(c)h_i(c)/(1-sp_i(c))$ and a direct channel computation gives
\[
 \mathcal D_i:=\mathbb E\operatorname{Var}(h_i\mid c)
 =\sum_cp_i(c)h_i(c)^2
 \left[A_i-\frac{s\nu(c)}{(1-sp_i(c))^2}\right],\quad
 A_i=\sum_c\frac{\nu(c)}{1-sp_i(c)}.
\]
Thus $\mathcal D_i\ge\kappa_{B_{\rm SG}}\sum_cp_i(c)h_i(c)^2$, where
$\kappa_{B_{\rm SG}}=1-B_{\rm SG}/(1-B_{\rm SG})^2\ge1-2B_{\rm SG}$.  Cauchy--Schwarz in the colour coordinate gives
\begin{equation}\label{eq:sg-channel-bound}
 \|Q\bar h\|^2\le \frac{s^2}{(1-B_{\rm SG})^2}
 \max_c\left[\nu(c)\sum_i p_i(c)\right]
 \sum_{i,c}p_i(c)h_i(c)^2.
\end{equation}
Let $u=sd/A_0$.  Since $A_0\ge s(d+\delta\Deg)$,
$u\le(1+\delta)^{-1}$.  Weighted AM--GM gives
\[
 Z\ge A_0\exp[-u/(1-B_{\rm SG})],\qquad
 a_c\sum_i sp_i(c)\prod_i(1-sp_i(c))\le e^{-1}.
\]
Consequently
\begin{equation}\label{eq:sg-star-occupancy}
 \frac{ds^2}{(1-B_{\rm SG})^2}\max_c\left[\nu(c)\sum_i p_i(c)\right]
 \le \frac{u}{(1-B_{\rm SG})^2}e^{u/(1-B_{\rm SG})-1}
 \le \frac{1}{1+3\delta/4}.
\end{equation}
(The last elementary inequality follows from $B_{\rm SG}\le\delta/[8(1+\delta)]$.)
Combining \eqref{eq:sg-channel-bound}--\eqref{eq:sg-star-occupancy} yields
\[
 d\|Q\bar h\|^2\le \frac{1}{(1+3\delta/4)\kappa_{B_{\rm SG}}}\|a\|^2
 \le \frac{1}{1+\delta/2}\|a\|^2,
\]
which, by \eqref{eq:sg-block-identity}, is exactly \eqref{eq:sg-additive-compression}.  All statements use only supported
coordinates, so the endpoint $x=0$ is included.
\end{proof}

\begin{lemma}[Unequal-incidence star inequality]
\label{lem:girth5-star-incidence}
Under the hypotheses of \Cref{lem:girth5-star-bounds}, put
\[
 \theta=\frac{1+\delta/4}{1+\delta/2}<1.
\]
If $\Deg\ge\Deg_0(\delta)$, then for every directed incidence choice
$0<\beta_i<2$ and every supported conditional star,
\begin{equation}\label{eq:sg-star-incidence}
 T_\beta:=Q+\tfrac12\{Q,W\}+D
 \succeq \frac{1}{2}Q-\frac{\theta}{2d}\sum_i\beta_i^2Q_i,
 \quad W=\sum_i\beta_iQ_i,\quad D=N^2-N,\ N=\sum_iQ_i,
\end{equation}
where $\theta$ is as above.
When $d=0$, the sum on the right is empty and the assertion means simply
$T_\beta=Q\succeq\frac12Q$; hence the displayed inequality is understood for
$d\ge1$.
\end{lemma}

\begin{proof}
We give the Schur-complement argument, including the constants.  On
$K_+$, put $K=\tfrac12H+\tfrac12(HW+WH)+D$.  Let
$\mathcal A$ and $\mathcal B$ be as above.  The assertion
$H\mathcal A=0$ uses the special form of $\mathcal A$, not merely the
fact that $H$ is a projection.  If
$a=h-\mathbb E[h\mid c]\in\mathcal A$ with
$h=\sum_i h_i(\sigma_i)$, then $Ph=h$ because $h$ is independent of the
centre spin.  Remove the global constant from $h$ (which does not change
$a$), and write $h=h_0+a$ with
$h_0=\mathbb E[h\mid c]\in K_0$, the $K_0$ and $K_+$ blocks of $Qh=0$
are
\[
 (I-P_{00})h_0=P_{0+}a,
 \qquad (I-P_{++})a=P_{+0}h_0.
\]
Consequently $Ca=P_{+0}(I-P_{00})^{-1}P_{0+}a=P_{+0}h_0$, and hence
$Ha=(I-P_{++}-C)a=0$.  Moreover,
$\mathcal A=\bigoplus_i\mathcal A_i$ with
$\mathcal A_i=\{h_i(\sigma_i)-\mathbb E[h_i\mid c]\}$, so $Q_i$ is the
identity on $\mathcal A_i$ and vanishes on $\mathcal A_j$ for $j\ne i$.
Thus $W$ is diagonal on $\mathcal A$; since it is self-adjoint, it also
preserves $\mathcal A^\perp=\mathcal B\oplus K_{\ge2}$.
On $\mathcal B$, writing $H_{11}=I-E$, $0\preceq E\preceq\eta I$, one has
\begin{equation}\label{eq:sg-singleton-block}
 K_{11}\succeq(1/2-5\eta/2)I.
\end{equation}
For $f\in K_{\ge2}$, completing the square gives
\[
 \frac{1}{2}\langle f,Hf\rangle+\Re\langle Hf,Wf\rangle
 \ge-\frac{1}{6}\|(W-I)f\|^2.
\]
Indeed, with $u=Hf$ and $b=(W-I)f$, the projection property
$H^*=H^2=H$ gives
\[
 \frac12\langle f,Hf\rangle+\Re\langle Hf,Wf\rangle
 =\frac32\|u\|^2+\Re\langle u,b\rangle
 =\frac32\left\|u+\frac13b\right\|^2-\frac16\|b\|^2,
\]
which proves the displayed bound.  On a conditional Hoeffding sector of
degree $n\ge2$, each $Q_i$ is either zero or the identity, so
$0\le W=\sum_i\beta_iQ_i\le2nI$.  Therefore
$(W-I)^2\preceq(2n-1)^2I\preceq\frac92n(n-1)I$, since the scalar
difference is
$\frac92n(n-1)-(2n-1)^2=\frac12(n-2)(n+1)\ge0$ for $n\ge2$.
Since $W$ and $D$ preserve these orthogonal sectors, this proves
$(W-I)^2\preceq\frac92D$ on all of $K_{\ge2}$.  Applying the
completed-square estimate to an arbitrary $f\in K_{\ge2}$ therefore gives
\[
 \langle f,Kf\rangle
 \ge \langle f,Df\rangle-\frac16\|(W-I)f\|^2
 \ge\frac14\langle f,Df\rangle.
\]
Thus $K_{\ge2,\ge2}\succeq D/4$, including all cross-sector terms of $H$.
The projection identity
$H_{1,\ge2}H_{\ge2,1}=H_{11}-H_{11}^2$ gives
$\|H_{1,\ge2}\|\le\sqrt\eta$.  Since $D\ge2I$ on the higher sectors,
\[
 \|K_{1,\ge2}D^{-1/2}\|
 \le \frac{1}{2}(3/\sqrt2+2\sqrt2)\sqrt\eta
 =\frac{7}{2\sqrt2}\sqrt\eta .
\]
To record the numerical margin, set
$X=K_{1,\ge2}D^{-1/2}$.  Since
$\|X\|^2\le(49/8)\eta$ and the higher block is at least $D/4$, retaining
$D/8$ leaves a Schur-complement loss at most
$8\|X\|^2\le49\eta$.  The singleton block left after this elimination
is therefore bounded below by
\[
 \left(\tfrac12-\tfrac52\eta\right)-\tfrac14-49\eta
 =\tfrac14-\tfrac{103}{2}\eta\ge0
 \qquad(\eta\le5/2048).
\]
Hence
\begin{equation}\label{eq:sg-coercivity}
 K\succeq \frac{1}{4}\Pi_{\mathcal B}+\frac{1}{8}D.
\end{equation}

On the other hand, for $f\in K_+$ write $f=a+y$, where
$a\in\mathcal A$ and $y\in K_+\ominus\mathcal A$.
Positivity of $C$ and Young's inequality imply
\begin{equation}\label{eq:sg-young}
 \langle Wf,CWf\rangle
 \le(1+\epsilon)\langle Wa,CWa\rangle
 +(1+\epsilon^{-1})r\|Wy\|^2.
\end{equation}
Because $W^2=\sum_i\beta_i^2Q_i$ on $\mathcal A$,
\eqref{eq:sg-additive-compression} bounds
$\langle Wa,CWa\rangle$ by $\theta_0\langle a,W^2a\rangle/d$.
Furthermore
$\|Wy\|^2\le4\|y_1\|^2+8\langle y,Dy\rangle$, where $y_1:=\Pi_{\mathcal B}y$ is the singleton component of $y$.
Take $\epsilon=\delta/4$.  The threshold gives
\[
 r\le R_\delta:=\frac{2(1+\delta)e^{2/\delta}}{\delta^3\Deg}
 \le\frac{\delta}{2048},\qquad
 \eta\le5/2048,
\]
so $(1+\epsilon^{-1})r\le1/32$.  We now verify the claimed combination
of \eqref{eq:sg-coercivity} and \eqref{eq:sg-young}.
To make the Schur cancellation explicit, set
$E_\beta=(\theta/d)\sum_i\beta_i^2Q_i$ and consider
$A:=T_\beta-Q/2+E_\beta/2$ on $K_0\oplus K_+$.  Its blocks are
\[
 \begin{aligned}
 A_{00}&=\tfrac12(I-P_{00}),
 &A_{0+}&=-\tfrac12P_{0+}(I+W),\\
 A_{++}&=\tfrac12(I-P_{++})+W-\tfrac12\{P_{++},W\}+D+\tfrac12E_\beta .
 \end{aligned}
\]
Since $A_{00}^{-1}=2(I-P_{00})^{-1}$, eliminating $K_0$ subtracts
$\tfrac12(I+W)C(I+W)$.  Using
$H=I-P_{++}-C$ and expanding the anticommutator, the resulting block is
\[
 \tfrac12H+\tfrac12\{H,W\}+D-\tfrac12WCW+\tfrac12E_\beta
 =K-\tfrac12WCW+\tfrac12E_\beta .
\]
\begin{equation}\label{eq:sg-schur-complement}
 K-\tfrac12WCW+\frac{\theta}{2d}\sum_i\beta_i^2Q_i\succeq0.
\end{equation}
Eliminating the degree-zero block (whose Schur complement is
$-\tfrac12WCW$) proves \eqref{eq:sg-star-incidence}.  This step is valid because the correction
vanishes on $K_0$.
\end{proof}

\begin{theorem}[Uniform spectral gap under the degree budget]
\label{thm:potts-gap-girth5}
Fix $0<\delta\le1$, $\Deg\ge\Deg_0(\delta)$ and
$q\ge(1+\delta)\Deg$.  Every finite Potts instance on a simple graph with
$\deg_H(v)+k_H(v)\le\Deg$ at each free vertex, after any normalized
pinning whose free residual graph has girth at least five, and for every
$x\in[0,1]$, has the nonnegative rate-one Glauber Laplacian
$\mathcal L:=\sum_v(I-P_v)$ satisfying
\begin{equation}\label{eq:sg-gap}
 \mathcal L^2\succeq\gamma_\delta \mathcal L,
 \qquad \gamma_\delta=\frac{\delta}{4(2+\delta)}.
\end{equation}
\end{theorem}

\begin{proof}
For each free vertex $v$ let $d_v$ be its current free degree and choose, on
an edge $vu$,
\[
 \beta_{vu}=\frac{2d_v}{d_v+d_u},\qquad
 W_v:=\sum_{u\sim v}\beta_{vu}Q_u,\qquad
 N_v:=\sum_{u\sim v}Q_u,\qquad
 T_v:=Q_v+\tfrac12\{Q_v,W_v\}+N_v^2-N_v.
\]
Girth at least five implies that every distance-two pair has a unique common
neighbour and that adjacent vertices have no common neighbour.  Therefore
\begin{equation}\label{eq:sg-global-decomposition}
 \mathcal L^2=\sum_vT_v+2\sum_{\{u,w\}:\,\operatorname{dist}_H(u,w)\ge3}Q_uQ_w.
\end{equation}
The last sum is positive semidefinite because nonadjacent heat-bath
projections commute.  For an edge $vu$, both endpoint degrees are
positive, and
\begin{equation}\label{eq:sg-incidence-sum}
 \frac{\beta_{vu}^2}{d_v}=\frac{1}{d_u}\bigl(1-(\beta_{vu}-1)^2\bigr),
 \qquad \sum_{v\sim u}\frac{\beta_{vu}^2}{d_v}\le1.
\end{equation}
For an isolated vertex $u$, the neighbour sum is empty.
Applying \eqref{eq:sg-star-incidence} in every conditional-star fibre and summing with \eqref{eq:sg-incidence-sum}
shows
\[
 \mathcal L^2\succeq \frac{1-\theta}{2}\mathcal L
 =\frac{\delta}{4(2+\delta)}\mathcal L.
\]
At $x>0$ the supported chain is irreducible. At $x=0$, remove from the
colour list of a free vertex $v$ the colours used by its $k_v$ fixed
neighbours, and denote the resulting list by $L_v$ and its size by
$\ell_v$.  Since at most
$k_v$ colours are removed,
\[
 \ell_v-d_v\ \ge\ q-k_v-d_v
 \ \ge\ q-\Deg
 \ \ge\ \delta\Deg\ \ge\ 2,
\]
where the last inequality is part of the choice of $\Deg_0(\delta)$.
Thus $\ell_v\ge d_v+2$ for every free vertex.  For completeness,
connectivity follows by induction on the number of free vertices.  Given
two proper list-colourings $\alpha,\beta$, choose a vertex $v$ and put
$c=\beta(v)$.  For every neighbour $u$ of $v$ with current colour
$\alpha(u)=c$, recolour $u$ (one at a time) to a colour in $L_u$ different
from $c$ and from all colours currently used by the neighbours of $u$;
at most $d_u+1$ colours are forbidden and $\ell_u\ge d_u+2$.  Now recolour
$v$ to $c$.  Keep $v$ fixed at $c$ and apply the induction hypothesis to
$H-v$ with the reduced lists
$L'_u=L_u\setminus\{c\}$ for neighbours of $v$ and $L'_u=L_u$ otherwise.
These lists satisfy $|L'_u|\ge\deg_{H-v}(u)+2$, and the current and target
restrictions are proper list-colourings.  The base case has one vertex.
Thus the proper-list-colouring support is connected, and
\eqref{eq:sg-gap} is the ordinary spectral-gap bound at the endpoint as well.
\end{proof}
